\section{Method}

\subsection{Overview}

DUCA-TAD is an in-forward acquisition plugin placed before a reusable TAD detector.
It has five parts: deploy-visible actionness sources, a temporal utility
selector, hard top-$K$ acquisition, an original-time sparse grid, and an audit
record. The detector architecture is reused, while sparse input coordinates and
metadata are made explicit.

\subsection{Deploy-Visible Acquisition Sources}

The selector receives cheap per-position descriptors. The zero-shot branch
computes $\pzs(t)$ from generic action/background prompts and a frozen
vision-language checkpoint. Prompt text, prompt hash, checkpoint hash, whether
THUMOS class names are used, and score-generation cost are recorded. The branch
is deploy-visible if it uses generic prompts and no target split labels. Other
branches may include motion magnitude, feature energy, low-resolution temporal
change, or a learned C3/PAction source. Learned C3/PAction scores are baselines
or ablations, not part of the zero-shot claim.

For an actionness-like source, the descriptor is
\begin{equation}
\begin{split}
    x_t = [&\pzs(t),\ \Delta \pzs(t),\ |\Delta \pzs(t)|,\
           \mathcal{H}(\pzs(t)),\\
           &\mathrm{motion}(t),\ \mathrm{energy}(t),\
           t/(\Tdense-1),\ \Mmask_t],
\end{split}
    \label{eq:source_features}
\end{equation}
where $\Mmask_t$ is the valid-position mask. A compact temporal encoder maps
$x_{1:\Tdense}$ to utility logits $g_t$, boundary logits $b_t$, risk logits
$r_t$, redundancy logits $\rho_t$, and optional budget logits.

\subsection{Hard In-Forward Selection}

The main contract fixes the requested budget to $\Kmain=384$. DUCA ranks valid
positions by detector-utility logits and hard-selects
\begin{equation}
    \Ssel = \mathrm{TopK}(g,\Kmain,\Mmask),
    \qquad |\Ssel|\le384.
    \label{eq:hard_topk}
\end{equation}
Optional center-radius or gap guards repair pathological holes using the same
in-forward scores. Dynamic $K$ is an ablation: it cannot replace the fixed-budget
main result, and every dynamic run is paired with a fixed run matched by average
selected count.

\begin{algorithm}[t]
\caption{DUCA in-forward acquisition for one detector forward}
\label{alg:duca_online}
\begin{algorithmic}[1]
\Require Dense candidate window $\VideoWin$, valid mask $\Mmask$, budget $\Kmain$
\State compute deploy-visible descriptors $x_{1:\Tdense}$
\State $(g,b,r,\rho) \gets \mathrm{DUCA}_{\thet}(x,\Mmask)$
\State $\Ssel \gets \mathrm{HardTopK}(g,\Kmain,\Mmask)$
\State $\Grid_{\Ssel} \gets \mathrm{BuildSparseGrid}(\Ssel,\Tdense,\Mmask)$
\State $Y \gets \Detector_\psi(\VideoWin_{\Ssel},\Grid_{\Ssel})$
\State write audit record $(\Ssel,\Grid_{\Ssel},\mathrm{source\ hashes},\mathrm{flags})$
\State \Return detector predictions $Y$ in original time
\end{algorithmic}
\end{algorithm}

\subsection{Detector-Utility Warm-Up}

The warm-up stage uses \cref{eq:utility_target} to initialize the selector.
Dense detector responsibility, teacher losses, boundary targets, and
hard-negative signals are available only on the training split. The warm-up loss
is
\begin{align}
    \mathcal{L}_{\mathrm{warm}}
    =&\ \lambda_u \mathrm{BCE}(g_t,u^{\star}_t)
      + \lambda_b \mathcal{L}_{\mathrm{bnd}}
      + \lambda_h \mathcal{L}_{\mathrm{hole}} \notag\\
     &+ \lambda_r \mathcal{L}_{\mathrm{red}}
      + \lambda_k \mathcal{L}_{\mathrm{budget}} .
    \label{eq:warmup_loss}
\end{align}
This stage teaches the selector that useful observations include boundaries,
hard negatives, and proposal-support regions, not only high-actionness interiors.

\subsection{Hard-Forward Surrogate Fine-Tuning}

Warm-up alone would leave the selector calibrated to a dense teacher but not to
the sparse detector regime. DUCA therefore fine-tunes with a hard sparse forward
pass: the detector consumes only $\VideoWin_{\Ssel}$, while the selector is
updated by explicit surrogate losses computed from train-only utility,
boundary, hole, and coverage targets on the same windows. Let $\pi_t=\sigma(g_t)$ and
\begin{equation}
    q_t
    =
    \mathrm{stopgrad}(\mathbf{1}[t\in\Ssel])
    + \pi_t - \mathrm{stopgrad}(\pi_t) .
    \label{eq:surrogate_mask}
\end{equation}
The packed hard subset is the only detector input. The detector loss trains the
detector under sparse observations, and the selector loss aligns the selected
set with detector-utility targets without assuming differentiability through
Top-$K$:
\begin{align}
    \mathcal{L}_{\mathrm{joint}}
    &=
    \mathcal{L}_{\mathrm{det}}
    \bigl(\Detector_\psi(\VideoWin_{\Ssel},\Grid_{\Ssel})\bigr)
    \notag\\
    &\quad
    + \lambda_s \mathcal{L}_{\mathrm{surrogate}}(q,u^\star)
    + \lambda_c \mathcal{L}_{\mathrm{cover}}(q).
    \label{eq:joint_loss}
\end{align}
The forward computation remains within the selected-observation budget. The
surrogate terms prevent the selector from learning only high-actionness
interiors. A custom straight-through bridge can be enabled as an implementation
variant only after a gradient test shows nonzero detector-loss gradients on the
selector parameters; without that proof, direct detector-loss-to-selector
gradient is not part of the main claim.

\subsection{Original-Time Sparse Grid}

The sparse grid in \cref{eq:sparse_grid,eq:grid_cells,eq:original_time_decode}
maps selected-axis predictions back to original time. If a head predicts offsets
around selected index $i$, the support cell $(\ell_i,r_i)$ and dense coordinate
$s_i$ define assignment, regression ranges, remapping, NMS, and mAP. This step
is mandatory for high-IoU evaluation because selected rank is not a metric time
axis.

\subsection{Teacher-Free Inference and Audit Ledger}

At inference, deploy-visible sources feed the DUCA selector, the selector emits
hard selected observations, and the detector predicts from that sparse input.
The audit ledger is written after selection and contains selected positions,
sparse-grid metadata, prompt/checkpoint hashes, source costs, and no-leak flags.
It is never read as a decision source. Recursive validators reject GT fields,
teacher payloads, raw detector predictions, prediction caches, oracle boundary
scores, and policy/checkpoint mismatches.
